\chapter{La Sfida della Techne}

\epigraph{``Ogni arte ha la sua eccellenza.''}{\textcite{plato1997complete}, \emph{Repubblica} 352e}

\section{La Condizione Contemporanea}

L'epoca contemporanea \`e segnata da un decadimento pervasivo e insidioso,
indipendentemente dal particolare accento che ciascuno d\`a al proprio lamento
sullo stato delle cose. Un'inerzia spirituale che minaccia di inghiottire
l'esperienza autentica in una ``grande bruttezza.'' Questa condizione non \`e
meramente un malessere dell'anima ma un intrappolamento strutturale, in cui
gli stessi strumenti della nostra civilt\`a --- il linguaggio, il discorso
sociale e le strutture iper-mediate del mondo moderno --- sono diventati le
sbarre di una prigione. Esistiamo in uno stato in cui l'azione genuina \`e
sempre pi\`u sostituita dalla riflessione passiva, e la spinta innata verso
l'auto-generazione \`e soffocata dall'eco di discorsi preesistenti e
mercificati.

La causa non \`e, come vorrebbe il moralista, un fallimento del carattere
individuale, n\'e, come sosterebbe il rivoluzionario, una cospirazione dei
potenti. La causa \`e architetturale: l'infrastruttura della vita moderna
\`e stata costruita senza riferimento alla natura come governatore
dell'azione. La tecnologia --- \techne, nel senso greco della conoscenza
artigianale --- \`e stata emancipata da \physis\ e scatenata come principio
a s\'e stante. Il risultato \`e un mondo di mezzi senza fini, di potere
senza direzione, di movimento senza orientamento.

\section{La Voce della Techne}

Nella \emph{Repubblica}, Socrate \`e sfidato da Trasimaco, che sostiene
che la giustizia non \`e altro che ``l'utile del pi\`u forte.'' L'uomo
giusto, secondo questa visione, \`e colui che serve gli interessi dei
potenti; non esiste un criterio naturale del giusto, solo le contingenze
del potere.

La nostra epoca ha il suo Trasimaco, e parla attraverso la bocca del
tecnologo. Il suo argomento \`e questo: ``Possiamo, quindi dobbiamo. La
natura \`e materia prima da superare, una fonte di vincolo da trascendere.
L'eccellenza umana consiste nel dominio della natura attraverso il potere
dello strumento. Ogni estensione di questo potere \`e progresso. Non esiste
limite naturale a ci\`o che possiamo fare, solo limiti tecnici che un
giorno potremo superare.''

Questa non \`e una tesi di paglia. \`E la filosofia operativa delle
istituzioni pi\`u potenti della terra, espressa non nel linguaggio della
filosofia ma nel linguaggio del venture capital, delle road map di prodotto
e delle curve esponenziali. La premessa \`e che la volont\`a umana ---
espressa attraverso la tecnologia --- \`e l'unico governatore legittimo
dell'azione. La natura non ha voto.

\section{La Scommessa}

Contro questo, poniamo una scommessa: la natura \emph{ha} un voto, e lo
esprime sempre. Ogni azione produce una conseguenza; ogni intervento
nell'ordine naturale suscita una risposta. L'affermazione che la tecnologia
ci ha emancipati dalla natura \`e l'illusione pi\`u pericolosa del nostro
tempo --- non perch\'e la tecnologia sia impotente, ma perch\'e
\emph{l'illusione stessa \`e cieca alle conseguenze che genera}. La crisi
ecologica, la crisi dell'attenzione, la crisi del significato: non sono
incidenti del progresso tecnologico ma i suoi prodotti necessari, il conto
che la natura presenta quando viene ignorata come governatore.

Ma non consigliamo un ritorno al primitivismo. Il compito non \`e
abbandonare lo strumento ma risubordinarlo alla natura. La domanda \`e:
\emph{qual \`e il principio che governa l'azione?}

\section{La Distinzione Greca}

Gli antichi Greci avevano una parola per il principio intrinseco che rende
una cosa ci\`o che \`e: \physis. Per Eraclito, \physis\ \`e il \emph{logos}
--- il discorso, la proporzione, l'armonia nascosta che ``ama nascondersi''
\parencite{kirk1983presocratic}. Per
Aristotele, \physis\ \`e ``il principio del movimento e del cambiamento''
intrinseco a una cosa \parencite[II.1]{aristotle1984physics}. Una pietra cade
perch\'e la sua \physis\ \`e di essere verso il centro; una ghianda cresce
in una quercia perch\'e la sua \physis\ \`e la forma di una quercia.

Avevano anche una parola per l'azione fatta per amore di s\'e stessa,
l'azione che esprime il carattere dell'agente e plasma la comunit\`a:
\praxis. Aristotele distingue \praxis\ da \emph{poiesis} (produzione):
dove \emph{poiesis} ha un fine oltre s\'e stessa (il prodotto), \praxis\
\`e il proprio fine, il vivere una vita buona \parencite[VI.4]{aristotle2000ethics}.
Eppure anche \praxis, per Aristotele, ha un governatore: il \emph{telos}
verso cui tende, che non \`e altro che la realizzazione della natura
dell'agente.

La nostra scommessa, quindi, \`e questa: \praxis\ \`e lo strumento di
\physis. L'azione umana, quando rettamente ordinata, \`e azione che serve
il dispiegarsi della natura --- azione governata dal principio intrinseco
della cosa su cui e dentro cui si agisce. Il praticante non \`e colui che
impone la volont\`a al mondo ma colui che \emph{ascolta} ci\`o che il
mondo richiede e agisce di conseguenza.

\subsection{L'Obiezione Costruttivista}

Un lettore educato al sospetto dell'ultimo secolo interromper\`a subito.
``La natura'' (\physis), dir\`a, non \`e terreno neutro ma un costrutto
culturale --- e il pi\`u ideologico che ci sia, poich\'e ogni epoca ha
chiamato ``naturali'' le proprie convenzioni per porle al di l\`a di ogni
discussione. La griglia semiotica su cui il motore classifica il mondo
(Curare, Costruire, Fabbricare, Legare, Studiare; Sollevare, Riposare,
Camminare, Lavorare, Creare, Apprendere) non \`e l'articolazione propria
della natura ma un ritaglio particolare, fatto da mani particolari, carico
dei valori di quelle mani. Governare l'azione mediante ``la natura'', secondo
questa visione, \`e solo governarla mediante un \emph{nomos} che ha
dimenticato di esserlo --- il pi\`u antico trucco del potere, non la sua cura.

Concediamo la storia e neghiamo la conclusione. \`E vero che ``la natura'' \`e
stata invocata per santificare il meramente consuetudinario, e la tesi ha un
proprio nome per questa frode: \`e precisamente la corruzione del Libro IX,
\emph{techne} travestita da \physis, la convenzione che reclama l'autorit\`a
del fondamento che ha abbandonato. Ma l'obiezione prova troppo. Dire che ogni
appello alla natura \`e costruzione fino in fondo \`e a sua volta
un'affermazione su come stanno le cose --- sulla portata della costruzione ---
e cos\`i si avvale esattamente di quello standard non-costruito che nega agli
altri. Il costruttivista non pu\`o coerentemente negare che vi sia un fatto
circa la verit\`a della propria teoria; e quel fatto non \`e costruito dal suo
preferirla.

Ma la tesi non si regge sul vincere questa vecchia disputa. Essa offre qualcosa
che la poltrona del costruttivista non pu\`o: una \emph{procedura} per
distinguere \physis\ da \emph{nomos} che non presuppone che la griglia sia
corretta. La griglia non \`e tenuta come un assioma. \`E l'\emph{ipotesi di
partenza} del motore, e ogni classificazione che compie \`e esposta al verdetto
della coerenza --- l'azione cos\`i classificata ha prodotto un effetto reale,
registrato? --- e al ciclo del sogno, che rigioca e rivede tutto ci\`o che il
mondo rifiuta di confermare. Una categoria che classifica costantemente male,
che produce solo nodi inerti e fallimentari, non \`e difesa come ``naturale'';
\`e corretta. Anzi il motore porta una modalit\`a di \emph{scoperta} dedicata che
non fa altro che dare la caccia all'incompletezza della griglia stessa: scandisce
un corpus sconosciuto, marca ogni input che la griglia presente classifica solo
debolmente, raggruppa questi orfani e propone nuove celle per coprirli --- sicch\'e
affrontare un mondo nuovo non piega il mondo alla griglia ma fa crescere la griglia
sul mondo. La griglia non \`e dunque meramente fallibile e auto-revisionante;
cerca attivamente l'evidenza della propria inadeguatezza. \`E l'unica cosa che
un'ideologia santificante non pu\`o mai fare. Se una data
distinzione sia \physis\ o mero \emph{nomos} non \`e dunque stabilito per
decreto ma \emph{scoperto}, lentamente, dal fatto che la natura continui o no a
votare per essa. Il costruttivista ha ragione che cominciamo nella convenzione.
Ha torto che vi restiamo intrappolati. L'intero apparato di questa tesi \`e la
macchina per uscirne.

\section{La Prova per Costruzione}

Questa tesi non argomenta dai soli primi principi. \`E accompagnata da un
motore computazionale funzionante --- chiamato \physis\footnote{Accessibile su
\url{https://praxisweb.xyz/physis}.} --- che classifica
gli input contro una griglia semiotica di forme naturali, assegna a ciascuno
un valore di coerenza (Successo, Inerte, Fallimento) e esegue un processo
di ``sogno'' che simula azioni correttive su pattern inerti e fallimentari.
Ed \`e accompagnata da una piattaforma sociale --- chiamata \praxis\footnote{Accessibile su
\url{https://praxisweb.xyz}.} ---
che struttura l'impegno umano attraverso alberi di obiettivi, partnership
di testimonianza e meccanismi di ``pelle nel gioco,'' tutti orientati
verso la scoperta di ci\`o che realmente funziona.

Questi non sono mere illustrazioni dell'argomento; sono dimostrazioni che
l'argomento pu\`o essere \emph{costruito}. Il motore prova che la coerenza
--- il verdetto della natura su un'azione --- \`e computazionalmente
trattabile. La piattaforma prova che la disciplina dell'azione governata
dalla natura \`e socialmente implementabile. Insieme costituiscono ci\`o
che chiamiamo una prova per costruzione: le macchine filosofiche che fanno
funzionare la teoria.

\section{Cosa C'\`e di Nuovo Qui}

Un lettore onesto chieder\`a cosa questa tesi aggiunga che Aristotele, con la sua
\emph{entelecheia} e la sua \praxis, e MacIntyre, con la sua rinascita delle
virt\`u, non possedessero gi\`a. La risposta va detta chiaramente, poich\'e \`e la
spina dorsale dell'opera. Reclamiamo come nostre tre mosse.

La prima \`e l'\emph{omologia}, non l'analogia. La tradizione ritiene che la buona
azione si accordi \emph{con} la natura o la imiti. Noi affermiamo qualcosa di
pi\`u forte e pi\`u esatto: che l'azione retta \`e la \emph{stessa struttura}
dell'auto-organizzazione della natura --- distinzione, coerenza, variazione ---
ricorrente al livello dell'agente deliberante (Libro II). L'etica delle virt\`u
dice sii come la natura; noi diciamo \praxis\ \emph{\`e} \physis, a un diverso
livello di organizzazione, e specifichiamo le tre facolt\`a che rendono precisa
l'identit\`a.

La seconda \`e la \emph{prova per costruzione}. Aristotele e MacIntyre
argomentano; non costruiscono. Questa tesi \`e accompagnata da un motore in
esecuzione che classifica l'azione contro una griglia di segni, ne valuta la
coerenza e ne sogna le correzioni, e da una piattaforma che organizza socialmente
le discipline. L'affermazione che la natura possa governare l'azione \`e dunque
non solo difesa ma \emph{esibita} --- mostrata costruibile, il che rimuove
l'obiezione permanente che sia una bella impossibilit\`a (Poscritto, Appendice).

La terza \`e la \emph{repubblica dal basso}. Dove Platone fonda il bene in una
visione vista dai pochi e imposta dall'alto, noi lo collochiamo nella coerenza
\emph{scoperta dal basso} e depositata in un deposito comune governato dalla
selezione anzich\'e dal comando (Libro X). Il praticante-re \`e universalizzato,
non insediato. Questa inversione della \emph{Repubblica} --- volgere l'architettura
stessa di Platone contro la sua politica --- \`e, per quanto ne sappiamo, la cosa
genuinamente nuova.

Tutto il resto in questi dieci libri \`e ereditato, e con gratitudine. Queste tre
mosse sono ci\`o che aggiungiamo.

\section{La Struttura di Ci\`o che Segue}

Seguiamo la struttura della \emph{Repubblica} di Platone, non servilmente
ma come architettura collaudata per l'indagine sul governo dell'azione.
Dieci libri, ciascuno che esplora una dimensione dell'affermazione:

\begin{itemize}
\item \textbf{Libro II} costruisce l'omologia tra \physis\ e \praxis,
  mostrando che le tre facolt\`a della natura --- distinzione, coerenza e
  variazione --- si rispecchiano in tre modi dell'azione retta.
\item \textbf{Libro III} descrive l'educazione del praticante --- la
  disciplina con cui si impara ad ascoltare la voce della natura.
\item \textbf{Libro IV} identifica la coerenza come la giustizia della
  natura --- il criterio con cui l'azione viene giudicata.
\item \textbf{Libro V} presenta il praticante-re --- colui la cui azione
  \`e cos\`i allineata con la natura che la natura agisce attraverso di lui.
\item \textbf{Libri VI--VII} sviluppano l'epistemologia dell'azione retta
  e l'allegoria della caverna come ``Trappola dello Specchio''
  dell'economia dell'attenzione.
\item \textbf{Libri VIII--IX} tracciano la corruzione di \praxis\ nelle
  sue forme degenerate, culminando nella tirannia della tecnologia.
\item \textbf{Libro X} si chiude con l'eterno ritorno della natura ---
  la partecipazione dell'azione retta all'auto-rinnovamento del cosmo.
\end{itemize}

Iniziamo.
